\documentclass[tikz,border=8pt]{standalone}
\usepackage{tikz}
\usepackage{amsmath}
\usepackage{amssymb}
\usepackage{pgfplots}
\pgfplotsset{compat=1.18}
\usepackage{ctex}
\usetikzlibrary{calc, positioning, arrows.meta}

\begin{document}
\begin{tikzpicture}

%% ===== 左图：唯一解 =====
\begin{scope}[xshift=0cm]
  \begin{axis}[
    name=ax1,
    width=5.2cm, height=5.2cm,
    xmin=-0.5, xmax=3.5,
    ymin=-0.5, ymax=3.5,
    xlabel={$x$}, ylabel={$y$},
    xlabel style={font=\scriptsize},
    ylabel style={font=\scriptsize},
    axis lines=center,
    xtick={1,2,3}, ytick={1,2,3},
    xticklabel style={font=\scriptsize},
    yticklabel style={font=\scriptsize},
    tick style={thin},
    grid=major, grid style={gray!25, thin},
    clip=false
  ]
    % 直线1: y = x - 0.5  (即 x - y = 0.5, 表示 2x-2y=1)
    \addplot[blue!70!black, thick, domain=-0.5:3.5, samples=2] {x - 0.5};
    % 直线2: y = -x + 3   (即 x + y = 3)
    \addplot[red!70!black, thick, dashed, domain=-0.5:3.5, samples=2] {-x + 3};
    % 交点 (1.75, 1.25)
    \addplot[black, only marks, mark=*, mark size=2pt] coordinates {(1.75,1.25)};
  \end{axis}
  % 交点标注 (outside axis)
  \node[font=\scriptsize, fill=white, draw=gray!50, thin, rounded corners=2pt,
        inner sep=3pt] at ([xshift=0.3cm, yshift=-0.9cm] ax1.north) {交点 $(x^*,y^*)$};
  % 方程标注
  \node[font=\scriptsize, blue!70!black] at ([xshift=-0.4cm, yshift=0.9cm] ax1.south east)
        {$\ell_1$};
  \node[font=\scriptsize, red!70!black]  at ([xshift=0.5cm, yshift=1.4cm] ax1.south west)
        {$\ell_2$};
  % 子图标题
  \node[font=\small\bfseries] at ([yshift=-0.55cm] ax1.south) {唯一解};
  \node[font=\scriptsize, fill=white, draw=gray!50, thin, rounded corners=2pt, inner sep=3pt]
        at ([yshift=-1.1cm] ax1.south) {两直线相交于一点};
\end{scope}

%% ===== 中图：无穷多解 =====
\begin{scope}[xshift=6cm]
  \begin{axis}[
    name=ax2,
    width=5.2cm, height=5.2cm,
    xmin=-0.5, xmax=3.5,
    ymin=-0.5, ymax=3.5,
    xlabel={$x$}, ylabel={$y$},
    xlabel style={font=\scriptsize},
    ylabel style={font=\scriptsize},
    axis lines=center,
    xtick={1,2,3}, ytick={1,2,3},
    xticklabel style={font=\scriptsize},
    yticklabel style={font=\scriptsize},
    tick style={thin},
    grid=major, grid style={gray!25, thin},
    clip=false
  ]
    % 两条重合直线: y = x - 0.5，用略偏移使视觉上显示为一条
    \addplot[blue!70!black, ultra thick, domain=-0.5:3.5, samples=2] {x - 0.5};
    \addplot[red!70!black, thick, dashed, domain=-0.5:3.5, samples=2] {x - 0.5};
  \end{axis}
  \node[font=\scriptsize, blue!70!black] at ([xshift=0.2cm, yshift=1.2cm] ax2.south east)
        {$\ell_1\equiv\ell_2$};
  % 无穷解标注
  \node[font=\small, fill=white, draw=gray!50, thin, rounded corners=2pt, inner sep=3pt]
        at ([xshift=0cm, yshift=-0.55cm] ax2.north) {$\infty$ 解};
  \node[font=\small\bfseries] at ([yshift=-0.55cm] ax2.south) {无穷多解};
  \node[font=\scriptsize, fill=white, draw=gray!50, thin, rounded corners=2pt, inner sep=3pt]
        at ([yshift=-1.1cm] ax2.south) {两直线完全重合};
\end{scope}

%% ===== 右图：无解 =====
\begin{scope}[xshift=12cm]
  \begin{axis}[
    name=ax3,
    width=5.2cm, height=5.2cm,
    xmin=-0.5, xmax=3.5,
    ymin=-0.5, ymax=3.5,
    xlabel={$x$}, ylabel={$y$},
    xlabel style={font=\scriptsize},
    ylabel style={font=\scriptsize},
    axis lines=center,
    xtick={1,2,3}, ytick={1,2,3},
    xticklabel style={font=\scriptsize},
    yticklabel style={font=\scriptsize},
    tick style={thin},
    grid=major, grid style={gray!25, thin},
    clip=false
  ]
    % 直线1: y = x - 0.5
    \addplot[blue!70!black, thick, domain=-0.5:3.5, samples=2] {x - 0.5};
    % 直线2（平行）: y = x + 1
    \addplot[red!70!black, thick, dashed, domain=-0.5:3.5, samples=2] {x + 1};
  \end{axis}
  \node[font=\scriptsize, blue!70!black] at ([xshift=0.2cm, yshift=1.2cm] ax3.south east)
        {$\ell_1$};
  \node[font=\scriptsize, red!70!black]  at ([xshift=0.2cm, yshift=1.8cm] ax3.south east)
        {$\ell_2$};
  % 无解标注
  \node[font=\small, fill=white, draw=gray!50, thin, rounded corners=2pt, inner sep=3pt]
        at ([xshift=0cm, yshift=-0.55cm] ax3.north) {无解};
  \node[font=\small\bfseries] at ([yshift=-0.55cm] ax3.south) {无解};
  \node[font=\scriptsize, fill=white, draw=gray!50, thin, rounded corners=2pt, inner sep=3pt]
        at ([yshift=-1.1cm] ax3.south) {两直线互相平行};
\end{scope}

%% 整体标题
\node[font=\large\bfseries] at (9.0, 5.6) {二元线性方程组的三种解的情形};

\end{tikzpicture}
\end{document}
